\chapter{CASUÍSTICA}
\label{cap:casuistica}

Durante o período de estágio supervisionado, realizado de 09 de março a
22 de maio de 2026 no setor de Clínica Médica de Equinos da EV-UFMG, foi
possível acompanhar 40 casos clínicos. A diversidade dos casos
acompanhados ao longo desse período refletiu a complexidade da rotina
hospitalar e contribuiu de maneira significativa para a consolidação dos
conhecimentos adquiridos durante a graduação. Os dados coletados,
previamente organizados em planilha, foram organizados e apresentados a
seguir em formato de gráficos e tabelas, contemplando a distribuição dos
animais por espécie, sexo, padrão racial, pelagem, origem, sistemas
orgânicos acometidos, afecções diagnosticadas e resolução de cada caso.

\section{Espécies}
\label{sec:especies}

Dos 40 animais acompanhados, 39 eram equinos (97,5\%) (\textit{Equus ferus
caballus}) e 1 era muar (2,5\%), híbrido resultante do cruzamento entre
equino e asinino (\textit{Equus africanus asinus}). Embora os muares não
representem uma espécie propriamente dita, optou-se por incluí-los como
categoria independente na análise da casuística, por se tratar de
animais com características e demandas clínicas particulares que
justificam sua distinção dentro do grupo atendido.

\begin{figure}[htb]
  \centering
  \caption{Distribuição dos animais atendidos quanto à espécie.}
  \label{graf:especies}
  \includegraphics[width=0.9\textwidth]{media/image5.png}
  \fonte{Da autora (2026).}
\end{figure}

\section{Sexo dos animais atendidos}
\label{sec:sexo}

Em relação ao sexo, observou-se predominância de machos, que
corresponderam a 28 pacientes (70\%), enquanto as fêmeas totalizaram 12
animais (30\%).

\begin{figure}[htb]
  \centering
  \caption{Distribuição dos animais atendidos quanto ao sexo.}
  \label{graf:sexo}
  \includegraphics[width=0.9\textwidth]{media/image2.png}
  \fonte{Da autora (2026).}
\end{figure}

\section{Idade dos animais atendidos}
\label{sec:idade}

Quanto à faixa etária, o grupo mais representativo foi o de animais
entre 6 e 15 anos, com 18 indivíduos (45\%), seguido pelos animais com
até 1 ano, que totalizaram 11 casos (27,5\%). Animais entre 2 e 5 anos
corresponderam a 5 casos (12,5\%) e aqueles com mais de 15 anos somaram
6 indivíduos (15\%). Para os animais resgatados, a estimativa de idade
foi realizada por meio de avaliação da arcada dentária.

\begin{figure}[htb]
  \centering
  \caption{Distribuição dos animais atendidos quanto à idade.}
  \label{graf:idade}
  \includegraphics[width=0.9\textwidth]{media/image16.png}
  \fonte{Da autora (2026).}
\end{figure}

\section{Padrão racial}
\label{sec:racial}

Em relação ao padrão racial, observou-se predominância de animais sem
raça definida (SRD), correspondendo a 20 indivíduos (50\%). Entre os
animais com raça definida, o Mangalarga Marchador apresentou maior
frequência, totalizando 6 casos (15\%), seguido pelo Friesian, com 5
animais (12,5\%), pelo Brasileiro de Hipismo, com 4 registros (10\%), e
pelo Quarto de Milha, com 3 casos (7,5\%). As raças Puro-Sangue Lusitano
e Puro-Sangue Argentino foram representadas por 1 animal cada (2,5\%).

\begin{figure}[htb]
  \centering
  \caption{Distribuição dos animais atendidos quanto ao padrão racial.}
  \label{graf:racial}
  \includegraphics[width=0.9\textwidth]{media/image4.png}
  \fonte{Da autora (2026).}
\end{figure}

\section{Pelagem}
\label{sec:pelagem}

Quanto à pelagem, a castanha foi a mais prevalente, correspondendo a 14
animais (35\%), seguida pela tordilha, com 7 casos (17,5\%), e pela
alazã, com 6 (15\%). As pelagens pampa e preta foram observadas em 5
animais cada (12,5\%), enquanto a baia esteve presente em 3 indivíduos
(7,5\%).

\begin{figure}[htb]
  \centering
  \caption{Distribuição dos animais atendidos quanto à pelagem.}
  \label{graf:pelagem}
  \includegraphics[width=0.9\textwidth]{media/image1.png}
  \fonte{Da autora (2026).}
\end{figure}

\section{Origem}
\label{sec:origem}

Em relação à origem, a maioria dos animais era proveniente de
proprietários externos, correspondendo a 26 casos (65\%). Animais
resgatados, em maior parte por ONGs, totalizaram 11 indivíduos (27,5\%).
3 animais (7,5\%) pertenciam à própria Escola de Veterinária da UFMG ou
à Fazenda Experimental Prof. Hélio Barbosa (FEPHB), localizada em
Igarapé/MG.

\begin{figure}[htb]
  \centering
  \caption{Distribuição dos animais atendidos quanto à origem.}
  \label{graf:origem}
  \includegraphics[width=0.9\textwidth]{media/image24.png}
  \fonte{Da autora (2026).}
\end{figure}

\section{Sistemas acometidos}
\label{sec:sistemas}

Dentre os sistemas orgânicos acometidos, o sistema locomotor foi o mais
prevalente, com 13 ocorrências, seguido pelo sistema digestório, com 11.
O sistema tegumentar somou 5 ocorrências, o sistema respiratório 4 e o
sistema circulatório três. Estado geral e neonatologia apresentaram 2
casos cada, enquanto sistema nervoso, sistema reprodutor, sistema
odontológico e sistema linfático foram representados por 1 caso cada. É
importante destacar que alguns animais apresentaram acometimento
simultâneo de mais de um sistema, razão pela qual o total de ocorrências
supera o número de pacientes acompanhados.

A predominância de afecções do sistema locomotor é condizente com o
perfil dos animais atendidos, uma vez que cavalos de esporte, de
trabalho e de resgate apresentam maior predisposição a esse tipo de
acometimento em razão das exigências físicas e das condições às quais
são submetidos. A alta frequência de síndromes cólicas no sistema
digestório, por sua vez, está em consonância com a literatura, que
aponta as afecções gastrointestinais como uma das principais causas de
atendimento emergencial em equinos.

\begin{figure}[htb]
  \centering
  \caption{Distribuição dos animais atendidos quanto ao sistema orgânico acometido.}
  \label{graf:sistemas}
  \includegraphics[width=0.9\textwidth]{media/image3.png}
  \fonte{Da autora (2026).}
\end{figure}

A seguir, são detalhados os diagnósticos identificados em cada sistema
orgânico acometido. A frequência relativa (f.r.) corresponde ao
percentual de ocorrência do diagnóstico em relação ao total de
ocorrências do respectivo sistema, ao passo que a frequência absoluta
(f.a.) representa o percentual em relação ao total geral de 44
ocorrências registradas durante o estágio.

\subsection{Sistema digestório}
\label{subsec:digestorio}

\begin{table}[htb]
  \centering
  \caption{Número de ocorrências (n), frequência relativa (f.r.) e frequência absoluta (f.a.) dos diferentes diagnósticos associados ao sistema digestório.}
  \label{tab:digestorio}
  \begin{tabular}{lccc}
    \toprule
    \textbf{Diagnóstico} & \textbf{n} & \textbf{f.r. (\%)} & \textbf{f.a. (\%)} \\
    \midrule
    Síndrome cólica      & 7 & 63,6 & 15,9 \\
    Salmonelose          & 2 & 18,2 & 4,5  \\
    Síndrome de má absorção & 1 & 9,1 & 2,3 \\
    Diarreia             & 1 & 9,1  & 2,3  \\
    \bottomrule
  \end{tabular}
  \fonte{Da autora (2026).}
\end{table}

Dentre os casos de sistema digestório, a síndrome cólica foi a afecção
mais frequente, abrangendo distintas apresentações clínicas, desde
quadros de resolução clínica até casos que evoluíram a óbito, eutanásia
ou encaminhamento cirúrgico. Os pacientes foram, em sua maioria,
submetidos ao protocolo padrão de atendimento, que inclui sondagem
nasogástrica, lavagem gástrica, fluidoterapia e protocolos analgésicos,
sendo as palpações transretais e os exames ultrassonográficos abdominais
ferramentas fundamentais para o suporte diagnóstico. Os dois casos de
salmonelose ocorreram em potros com quadros sistêmicos graves, ambos
associados à coinfecção respiratória por \textit{Rhodococcus equi}, o que
limitou as opções terapêuticas e culminou em desfechos desfavoráveis. Já
o caso de síndrome de má absorção encontrava-se em investigação ao
término do estágio, ao passo que o quadro de diarreia evoluiu
favoravelmente, com alta clínica após estabilização hidroeletrolítica.

\subsection{Sistema locomotor}
\label{subsec:locomotor}

\begin{table}[htb]
  \centering
  \caption{Número de ocorrências (n), frequência relativa (f.r.) e frequência absoluta (f.a.) dos diferentes diagnósticos associados ao sistema locomotor.}
  \label{tab:locomotor}
  \begin{tabular}{lccc}
    \toprule
    \textbf{Diagnóstico} & \textbf{n} & \textbf{f.r. (\%)} & \textbf{f.a. (\%)} \\
    \midrule
    Laminite                               & 3 & 23,1 & 6,8 \\
    Lesões em TFDP                         & 2 & 15,4 & 4,5 \\
    Osso navicular bipartido               & 1 & 7,7  & 2,3 \\
    Fratura de processo espinhoso          & 1 & 7,7  & 2,3 \\
    Osteoartrite                           & 1 & 7,7  & 2,3 \\
    Contratura tendínea                    & 1 & 7,7  & 2,3 \\
    Pododermatite séptica                  & 1 & 7,7  & 2,3 \\
    Osteófito no terço distal do 3º metacarpo & 1 & 7,7 & 2,3 \\
    Lesão de TFDS                          & 1 & 7,7  & 2,3 \\
    Osteodistrofia fibrosa                 & 1 & 7,7  & 2,3 \\
    \bottomrule
  \end{tabular}
  \fonte{Da autora (2026).}
\end{table}

O sistema locomotor foi o mais prevalente entre os sistemas orgânicos
acometidos, refletindo o perfil dos animais atendidos no setor, em sua
maioria cavalos de esporte, de trabalho e de resgate. A laminite figurou
como a principal afecção, com manifestações de origens distintas. As
lesões em tendões flexores também foram recorrentes, sendo conduzidas
com base em avaliação ultrassonográfica. Os demais casos apresentaram
apenas uma ocorrência, mas evidenciam a diversidade de afecções
locomotoras acompanhadas, abrangendo desde alterações ortopédicas, como
osteoartrite e osteodistrofia fibrosa, até casos de encaminhamento
cirúrgico, como a fratura de processo espinhoso e a contratura tendínea
associada à lesão do ligamento acessório do tendão flexor digital
profundo.

\subsection{Sistema tegumentar}
\label{subsec:tegumentar}

\begin{table}[htb]
  \centering
  \caption{Número de ocorrências (n), frequência relativa (f.r.) e frequência absoluta (f.a.) dos diferentes diagnósticos associados ao sistema tegumentar.}
  \label{tab:tegumentar}
  \begin{tabular}{lccc}
    \toprule
    \textbf{Diagnóstico} & \textbf{n} & \textbf{f.r. (\%)} & \textbf{f.a. (\%)} \\
    \midrule
    Manejo de feridas & 3 & 60,0 & 6,8 \\
    Habronemose       & 1 & 20,0 & 2,3 \\
    Corpo estranho    & 1 & 20,0 & 2,3 \\
    \bottomrule
  \end{tabular}
  \fonte{Da autora (2026).}
\end{table}

No sistema tegumentar, predominou a necessidade de manejo de feridas,
cuja conduta envolveu, principalmente, limpeza com solução fisiológica e
clorexidina degermante, desbridamento, aplicação de pomadas e óleos
cicatrizantes, unguento e proteção com curativos ou bandagens, quando
aplicável. O caso de habronemose, observado em animal pertencente ao
próprio Hospital Veterinário, foi conduzido clinicamente com terapia
tópica e antiparasitária, permanecendo em tratamento ao final do
estágio. Já o caso de corpo estranho demandou encaminhamento ao setor de
Clínica Cirúrgica de Grandes Animais para tratamento definitivo da
lesão.

\subsection{Sistema respiratório}
\label{subsec:respiratorio}

\begin{table}[htb]
  \centering
  \caption{Número de ocorrências (n), frequência relativa (f.r.) e frequência absoluta (f.a.) dos diferentes diagnósticos associados ao sistema respiratório.}
  \label{tab:respiratorio}
  \begin{tabular}{lccc}
    \toprule
    \textbf{Diagnóstico} & \textbf{n} & \textbf{f.r. (\%)} & \textbf{f.a. (\%)} \\
    \midrule
    Rodococose        & 2 & 50,0 & 4,5 \\
    Hemiplegia laríngea & 2 & 50,0 & 4,5 \\
    \bottomrule
  \end{tabular}
  \fonte{Da autora (2026).}
\end{table}

As afecções respiratórias acompanhadas distribuíram-se igualmente entre
rodococose e hemiplegia laríngea. Os dois casos de rodococose ocorreram
em potros jovens, com quadro de broncopneumonia multifocal associada à
formação de abscessos pulmonares, sendo a confirmação diagnóstica obtida
por meio de exames ultrassonográficos, radiográficos e por achados
anatomopatológicos post-mortem. Ambos os pacientes apresentaram desfecho
desfavorável, evoluindo para óbito e eutanásia. Os casos de hemiplegia
laríngea, por sua vez, foram diagnosticados em equinos jovens da raça
Friesian encaminhados ao setor exclusivamente para realização do
procedimento diagnóstico, sendo a confirmação obtida por meio de
avaliação endoscópica estática e dinâmica das vias aéreas superiores.

\subsection{Sistema circulatório}
\label{subsec:circulatorio}

\begin{table}[htb]
  \centering
  \caption{Número de ocorrências (n), frequência relativa (f.r.) e frequência absoluta (f.a.) dos diferentes diagnósticos associados ao sistema circulatório.}
  \label{tab:circulatorio}
  \begin{tabular}{lccc}
    \toprule
    \textbf{Diagnóstico} & \textbf{n} & \textbf{f.r. (\%)} & \textbf{f.a. (\%)} \\
    \midrule
    Perfuração de jugular & 1 & 33,3 & 2,3 \\
    Tromboflebite         & 1 & 33,3 & 2,3 \\
    Hemoparasitose        & 1 & 33,3 & 2,3 \\
    \bottomrule
  \end{tabular}
  \fonte{Da autora (2026).}
\end{table}

No sistema circulatório, foram acompanhados três casos distintos. O caso
de perfuração de jugular evoluiu de maneira desfavorável, culminando em
eutanásia. A tromboflebite, complicação secundária à manipulação venosa
prévia, ocorreu de forma concomitante a um quadro de síndrome cólica,
sendo ambas as condições manejadas simultaneamente durante o período de
internação. O caso de hemoparasitose foi tratado conforme protocolo
preconizado para a doença, com boa resposta clínica e alta hospitalar.

\subsection{Estado geral}
\label{subsec:estado-geral}

\begin{table}[htb]
  \centering
  \caption{Número de ocorrências (n), frequência relativa (f.r.) e frequência absoluta (f.a.) dos diferentes diagnósticos associados ao estado geral.}
  \label{tab:estado-geral}
  \begin{tabular}{lccc}
    \toprule
    \textbf{Diagnóstico} & \textbf{n} & \textbf{f.r. (\%)} & \textbf{f.a. (\%)} \\
    \midrule
    Desnutrição & 2 & 100,0 & 4,5 \\
    \bottomrule
  \end{tabular}
  \fonte{Da autora (2026).}
\end{table}

Os casos relacionados ao estado geral envolveram animais provenientes de
resgate, com histórico de manejo inadequado, negligência e suspeita de
maus-tratos, evidenciados por quadro de privação alimentar prolongada e
debilidade generalizada. A conduta clínica priorizou a reintrodução
cuidadosa da alimentação, com suporte nutricional adequado e
monitoramento de parâmetros laboratoriais. Ambos os pacientes
apresentaram evolução clínica favorável e receberam alta hospitalar.

\subsection{Neonatologia}
\label{subsec:neonatologia}

\begin{table}[htb]
  \centering
  \caption{Número de ocorrências (n), frequência relativa (f.r.) e frequência absoluta (f.a.) dos diferentes diagnósticos associados aos neonatos.}
  \label{tab:neonatologia}
  \begin{tabular}{lccc}
    \toprule
    \textbf{Diagnóstico} & \textbf{n} & \textbf{f.r. (\%)} & \textbf{f.a. (\%)} \\
    \midrule
    Parto assistido  & 1 & 50,0 & 2,3 \\
    Cuidados neonatais & 1 & 50,0 & 2,3 \\
    \bottomrule
  \end{tabular}
  \fonte{Da autora (2026).}
\end{table}

Os casos de neonatologia corresponderam a um único evento obstétrico no
qual uma égua receptora foi encaminhada ao Hospital Veterinário para
parto assistido devido à suspeita de desproporção feto-pélvica. Após o
parto, que transcorreu de forma eutócica, foram realizados os cuidados
neonatais, incluindo cura de umbigo, avaliação da transferência de
imunidade passiva por meio de dosagem de imunoglobulinas (IgG check) e
suporte clínico à égua e ao neonato. Ambos evoluíram favoravelmente e
receberam alta hospitalar.

\subsection{Demais sistemas orgânicos}
\label{subsec:demais-sistemas}

\begin{table}[htb]
  \centering
  \caption{Número de ocorrências (n), frequência relativa (f.r.) e frequência absoluta (f.a.) dos diferentes diagnósticos associados aos demais sistemas (nervoso, reprodutor, odontológico e linfático).}
  \label{tab:demais-sistemas}
  \begin{tabular}{lccc}
    \toprule
    \textbf{Diagnóstico} & \textbf{n} & \textbf{f.r. (\%)} & \textbf{f.a. (\%)} \\
    \midrule
    Sinais neurológicos inespecíficos & 1 & 25,0 & 2,3 \\
    Parafimose                        & 1 & 25,0 & 2,3 \\
    Fratura dentária                  & 1 & 25,0 & 2,3 \\
    Linfoma                           & 1 & 25,0 & 2,3 \\
    \bottomrule
  \end{tabular}
  \fonte{Da autora (2026).}
\end{table}

Foram acompanhados ainda quatro casos relacionados a sistemas orgânicos
de menor prevalência na casuística do período. O caso neurológico, em
que se observaram sinais clínicos inespecíficos, evoluiu favoravelmente,
com remissão parcial do quadro e alta clínica. A parafimose, observada
em equino macho, foi encaminhada ao setor de Clínica Cirúrgica de
Grandes Animais para resolução definitiva. O quadro odontológico
envolveu fratura dentária em paciente resgatado, sendo conduzido por
meio de atendimento odontológico especializado com subsequente alta
hospitalar. Por fim, o caso de linfoma, já em estágio avançado no
momento do atendimento, teve desfecho desfavorável.

\section{Desfecho}
\label{sec:desfecho}

Quanto ao desfecho, a maioria dos animais recebeu alta médica,
totalizando 22 casos (55\%). Sete animais (17,5\%) permaneciam em
tratamento ao término do período de estágio e 6 (15\%) foram
encaminhados ao setor de Clínica Cirúrgica de Grandes Animais. Três
animais (7,5\%) foram submetidos à eutanásia e 2 (5\%) evoluíram a
óbito.

\begin{figure}[htb]
  \centering
  \caption{Distribuição dos animais atendidos quanto ao desfecho do caso.}
  \label{graf:desfecho}
  \includegraphics[width=0.9\textwidth]{media/image9.png}
  \fonte{Da autora (2026).}
\end{figure}
